\chapter{Grafikus felület specifikációja}

\section{A grafikus interfész}
\diagram{docs/img/mainmenu}{Főmenü}{15cm}
\diagram{docs/img/ingame}{Játékban}{15cm}

\section{A grafikus rendszer architektúrája}
\subsection{A felület működési elve}
A grafikus felület a meglévő üzleti logikát (\texttt{motor}, \texttt{jarmuvek}, \texttt{jatekosok}, \texttt{takaritofejek}, \texttt{terkep} csomagok) változtatás nélkül használja fel — a modell rétegben a GUI miatt csak additív bővítések történtek (új interfész, új getterek, új mezők), egyetlen meglévő metódus viselkedése sem módosult. Így a korábbi parancssoros prototípus ugyanúgy működik, mint korábban.

A három réteg felelőssége

A felület tiszta MVC elrendezést követ, három különálló csomagra bontva:
\begin{itemize}
    \item Modell (\texttt{motor}, \texttt{jarmuvek}, \texttt{jatekosok}, \texttt{takaritofejek}, \texttt{terkep}) — az állapot és a játékszabályok.
    \item Nézet (\texttt{gui.view}) — a Swing-panelek (\texttt{MapPanel}, \texttt{HudPanel}, \texttt{ContextPanel}, \texttt{EventLogPanel}), amelyek a modell aktuális állapotát rajzolják. A nézet csak olvas a modellből, soha, sehol nem módosítja közvetlenül.
    \item Vezérlő (\texttt{gui.controller}) — a \texttt{MapController} és az \texttt{ActionController}. Ezek a felhasználói inputot fordítják le model-hívásokká (pl. \texttt{Hokotro.takaritSav}, \texttt{Jarmu.moveToNode}, \texttt{GameActions.endRound}).
\end{itemize}

A megjelenítés alapelve: kevert (push + pull)

Push oldal — eseményalapú értesítés. A motor csomagba felvettünk egy \texttt{GameStateListener} interfészt egyetlen \texttt{onStateChanged(GameState state)} metódussal. A \texttt{GameState} egy \texttt{addListener} / \texttt{fireStateChanged} páros segítségével tartja számon a feliratkozott listenereket. Amikor a játékállapot ténylegesen változik (jármű kijelölése, lépés, takarítás, kör vége, vásárlás), egy magasabb szintű művelet végén — jellemzően az \texttt{ActionController}-ben vagy a \texttt{GameActions.endRound}-ban — a vezérlő meghívja a \texttt{state.fireStateChanged()}-et. A \texttt{GameWindow} minden view-panelt feliratkoztat listenerként, tehát az értesítés minden megjelenítendő komponensnek kézbesítődik.

Pull oldal — közvetlen lekérdezés rajzoláskor. A view-panelek \texttt{onStateChanged} implementációja nem ad át adatot magával az eseménnyel; egyszerűen kiváltja a saját frissítését (\texttt{repaint()} a \texttt{MapPanel}-en, vagy a label-ek és gombok újraépítése a \texttt{ContextPanel}-ben). A tényleges állapotot a panel a saját rajzolási vagy frissítési ciklusában közvetlenül olvassa le a modellből: a \texttt{MapPanel.paintComponent} például a \texttt{state.getAllUtak()}, \texttt{state.getAllEntities()}, \texttt{Sav.ho}, \texttt{Sav.ice} értékeit kérdezi le, a \texttt{ContextPanel} pedig a kiválasztott entitást a \texttt{state.selected()}-en keresztül. Így a megjelenítés mindig az aktuális, friss állapotot tükrözi anélkül, hogy a model-eseménynek bármilyen tartalmat kéne hordoznia.

\subsection{A felület osztály-struktúrája}
\diagram{docs/img/uml_class/view_vegleges}{Osztálydiagram}{16cm}

\section{A grafikus objektumok felsorolása}
\subsection{Új osztályok}

%------------ÚJ OSZTÁLYOK----------------

\subsubsection{GameStateListener}
\begin{class-template-responsibility}
    Push-event interfész a model $\rightarrow$ view kommunikáció megvalósításához. Lehetővé teszi, hogy a modell értesítse a feliratkozott nézeteket (view-kat) az állapotváltozásról, így azok frissíthetik a megjelenített adatokat.
\end{class-template-responsibility}
\begin{class-template-interface}
    -
\end{class-template-interface}
\begin{class-template-baseclass}
    -
\end{class-template-baseclass}
\begin{class-template-attribute}
    -
\end{class-template-attribute}
\begin{class-template-method}
    \classitem{+onStateChanged() : void}{A modell hívja meg, amikor az állapota megváltozik, jelezve a nézetnek a grafikus frissítés szükségességét.}
\end{class-template-method} 

\subsubsection{Main}
\begin{class-template-responsibility}
    A GUI alapú alkalmazás belépési pontja. Feladata a rendszerszintű megjelenés (Look & Feel) beállítása és a főmenü elindítása az Event Dispatch Thread-en (EDT).
\end{class-template-responsibility}
\begin{class-template-interface}
    -
\end{class-template-interface}
\begin{class-template-baseclass}
    -
\end{class-template-baseclass}
\begin{class-template-attribute}
    -
\end{class-template-attribute}
\begin{class-template-method}
    \classitem{+main(args : String[]) : void}{A program belépési metódusa.}
\end{class-template-method} 

\subsubsection{MainMenu}
\begin{class-template-responsibility}
    A játék indítóképernyője. Lehetővé teszi a nehézségi szint és a kezdő hókotrófej kiválasztását, majd elindítja a játékablakot.
\end{class-template-responsibility}
\begin{class-template-interface}
    -
\end{class-template-interface}
\begin{class-template-baseclass}
    JFrame
\end{class-template-baseclass}
\begin{class-template-attribute}
    \classitem{-difficultyCombo : JComboBox}{Nehézség választó (Könnyű, Közepes, Nehéz).}
    \classitem{-headRadioGroup : ButtonGroup}{Rádiógomb csoport a kezdő fej kiválasztásához (Söprő/Jégtörő).}
    \classitem{-startButton : JButton}{A játék indítását kiváltó gomb.}
\end{class-template-attribute}\begin{class-template-method}
    \classitem{+MainMenu()}{Konstruktor, felépíti a menü grafikus felületét.}
    \classitem{-onStartGame() : void}{Eseménykezelő, amely létrehozza a GameWindow-t a választott paraméterekkel, majd bezárja a főmenüt.}
\end{class-template-method}

\subsubsection{GameWindow}
\begin{class-template-responsibility}
    A játék fő ablaka, amely összefogja a különböző nézeteket és kontrollereket. Összeállítja a játék kezdeti állapotát és kezeli a panelek elrendezését.
\end{class-template-responsibility}
\begin{class-template-interface}
    -
\end{class-template-interface}
\begin{class-template-baseclass}
    JFrame
\end{class-template-baseclass}
\begin{class-template-attribute}
    \classitem{-state : GameState}{A játék aktuális modell-állapota.}
    \classitem{-mapPanel : MapPanel}{A térkép megjelenítő panel.}
    \classitem{-hudPanel : HudPanel}{A felső információs sáv.}
    \classitem{-contextPanel : ContextPanel}{A jobb oldali interakciós panel.}
    \classitem{-eventLogPanel : EventLogPanel}{Az alsó eseménynapló panel.}
\end{class-template-attribute}
\begin{class-template-method}
    \classitem{+GameWindow(diff : Difficulty, head : HeadType)}{Létrehozza az ablakot, inicializálja a modellt és a nézeteket, majd feliratkoztatja a paneleket a modell eseményeire.}
\end{class-template-method}

\subsubsection{GameInitializer}
\begin{class-template-responsibility}
    Statikus segédosztály, amely a specifikáció szerinti várostérképet és a kezdő objektumokat (csomópontok, utak, járművek, NPC-k) hozza létre a memóriában.
\end{class-template-responsibility}
\begin{class-template-interface}
    -
\end{class-template-interface}
\begin{class-template-baseclass}
    -
\end{class-template-baseclass}
\begin{class-template-attribute}
    -
\end{class-template-attribute}
\begin{class-template-method}
    \classitem{+createCity(state : GameState) : void}{Feltölti a kapott állapotot a hardcoded adatokkal: Telephely, Főtér, Vasútállomás, stb., valamint a játékosok és NPC-k kezdőhelyzeteivel.}
\end{class-template-method}

\subsubsection{GameLayout}
\begin{class-template-responsibility}
    A térkép statikus elhelyezkedéséért felelős segédosztály. Tárolja a koordinátákat és a megjelenítéshez szükséges konstansokat.
\end{class-template-responsibility}
\begin{class-template-interface}
    -
\end{class-template-interface}
\begin{class-template-baseclass}
    -
\end{class-template-baseclass}
\begin{class-template-attribute}
    \classitem{+NODE_RADIUS : int}{Csomópontok rajzolási sugara.}
    \classitem{+LANE_WIDTH : int}{Sávok szélessége pixelben.}
    \classitem{-nodePositions : Map<String, Point>}{Csomópont azonosítóhoz rendelt koordináták.}
    \classitem{-displayNames : Map<String, String>}{Belső azonosítók magyar megfelelői.}
\end{class-template-attribute}
\begin{class-template-method}
    \classitem{+getPosition(nodeId : String) : Point}{Visszaadja egy adott csomópont koordinátáját.}
    \classitem{+getDisplayName(id : String) : String}{Visszaadja az azonosító magyar nevét.}
\end{class-template-method}

\subsubsection{MapPanel}
\begin{class-template-responsibility}
    A város térképének és a rajta lévő egységeknek a vizuális megjelenítése. Kirajzolja a csomópontokat, az utak állapotát és a járművek aktuális pozícióját.
\end{class-template-responsibility}
\begin{class-template-interface}
    GameStateListener
\end{class-template-interface}
\begin{class-template-baseclass}
    JPanel
\end{class-template-baseclass}
\begin{class-template-attribute}
    \classitem{-state : GameState}{Referencia a modellre az adatok lekéréséhez.}
\end{class-template-attribute}
\begin{class-template-method}
    \classitem{#paintComponent(g : Graphics) : void}{Kirajzolja a térképet: sárga/zöld/szürke csomópontok, utak a színkódolt állapotuk szerint, és a járművek ikonjai.}
    \classitem{+vehicleAt(p : Point) : String}{Visszaadja a megadott koordinátán található jármű nevét, ha nincs ott semmi, null-t.}
    \classitem{+onStateChanged() : void}{A modell jelzésére újrarajzolja a panelt (repaint).}
\end{class-template-method}

\subsubsection{HudPanel}
\begin{class-template-responsibility}
    A játék felső állapotsávja, amely az általános információkat (kör, aktív játékos, vagyon, balesetek) jeleníti meg.
\end{class-template-responsibility}
\begin{class-template-interface}
    GameStateListener
\end{class-template-interface}
\begin{class-template-baseclass}
    JPanel
\end{class-template-baseclass}
\begin{class-template-attribute}
    \classitem{-roundLabel : JLabel}{Kör sorszámát mutató felirat.}
    \classitem{-wealthLabel : JLabel}{A játékosok vagyonát mutató felirat.}
    \classitem{-doneButton : JButton}{A kört befejező "Kész vagyok" gomb.}
\end{class-template-attribute}
\begin{class-template-method}
    \classitem{+onStateChanged() : void}{Frissíti a feliratokat és a gomb szövegét az aktív játékos neve alapján.}
\end{class-template-method}

\subsubsection{ContextPanel}
\begin{class-template-responsibility}
    A kiválasztott jármű részleteit és a rajta végezhető akciókat (mozgás, takarítás, telephelyi műveletek) megjelenítő panel.
\end{class-template-responsibility}
\begin{class-template-interface}
    GameStateListener
\end{class-template-interface}
\begin{class-template-baseclass}
    JPanel
\end{class-template-baseclass}
\begin{class-template-attribute}
    \classitem{-actionContainer : JPanel}{Dinamikusan változó gombok tárolója.}
    \classitem{-infoArea : JTextArea}{Jármű adatainak (üzemanyag, töltet) megjelenítése.}
\end{class-template-attribute}
\begin{class-template-method}
    \classitem{+onStateChanged() : void}{A kijelölt jármű típusa és helyzete alapján újragenerálja az akciógombokat.}
\end{class-template-method}

\subsubsection{EventLogPanel}
\begin{class-template-responsibility}
    A játék eseményeinek szöveges naplózása és a jelmagyarázat megjelenítése.
\end{class-template-responsibility}
\begin{class-template-interface}
    GameStateListener
\end{class-template-interface}
\begin{class-template-baseclass}
    JPanel
\end{class-template-baseclass}
\begin{class-template-attribute}
    \classitem{-logArea : JTextArea}{Az eseményeket listázó szövegdoboz.}
\end{class-template-attribute}
\begin{class-template-method}
    \classitem{+onStateChanged() : void}{Lekéri a modell legfrissebb eseményeit és megjeleníti azokat.}
\end{class-template-method}

\subsubsection{MapController}
\begin{class-template-responsibility}
    A térképen történő egérkattintások kezelése, a járművek kijelölésének és deszelektálásának vezérlése.
\end{class-template-responsibility}
\begin{class-template-interface}
    MouseListener
\end{class-template-interface}
\begin{class-template-baseclass}
    MouseAdapter
\end{class-template-baseclass}
\begin{class-template-attribute}
    \classitem{-mapPanel : MapPanel}{A megfigyelt térkép panel.}
    \classitem{-state : GameState}{A modell, ahol a kijelölést állítja.}
\end{class-template-attribute}
\begin{class-template-method}
    \classitem{+mouseClicked(e : MouseEvent) : void}{Kattintáskor a vehicleAt segédmetódussal lekérdezi, van-e ott jármű, és frissíti a modellben a kijelölt objektumot.}
\end{class-template-method}

\subsubsection{ActionController}
\begin{class-template-responsibility}
    A felhasználói felületről érkező gombnyomásokat fordítja le konkrét modell-műveletekre, kezeli a körök közötti váltást és a játék végének logikáját.
\end{class-template-responsibility}
\begin{class-template-interface}
    -
\end{class-template-interface}
\begin{class-template-baseclass}
    -
\end{class-template-baseclass}
\begin{class-template-attribute}
    \classitem{-state : GameState}{A játék modell-állapota.}
\end{class-template-attribute}
\begin{class-template-method}
    \classitem{+onMoveTo(roadId : String) : void}{A kijelölt jármű mozgatása a kiválasztott útra.}
    \classitem{+onTakarit() : void}{A takarítási akció elindítása a felszerelt fejnek megfelelően.}
    \classitem{+onBuyHokotro() : void}{Új hókotró vásárlása a telephelyen.}
    \classitem{+onPlayerDone() : void}{Az aktív játékos körének lezárása, váltás a következő játékosra, vagy NPC körök futtatása.}
    \classitem{-ensureControllable() : boolean}{Ellenőrzi, hogy a kijelölt jármű az aktív játékosé-e.}
    \classitem{-showGameOver(winner : String) : void}{Megjeleníti a játék vége dialógust a győztessel.}
\end{class-template-method}  

\subsection{Megváltozott osztályok}
\subsubsection{GameState}
\begin{class-template-responsibility}
    A játék teljes állapotát reprezentáló osztály, amely a módosítások során kibővült a nézetek értesítéséért felelős push-mechanizmussal (Listener minta), BFS alapú útvonalkereséssel az NPC-k számára, valamint a Pass'N'Play körökre bontott vezérléshez szükséges állapotokkal és a UI eseménynapló kezelésével.
\end{class-template-responsibility}
\begin{class-template-interface}
    -
\end{class-template-interface}
\begin{class-template-baseclass}
    -
\end{class-template-baseclass}
\begin{class-template-attribute}
    \classitem{+currentRound}{Körszámláló, értéke 1-től indul.}
    \classitem{+activePlayerName}{A Pass'N'Play módban éppen aktív játékos neve.}
    \classitem{+gameOverReason}{Magyar nyelvű szöveges üzenet a játék végének okáról.}
    \classitem{+running}{Jelzi, hogy a játék fut-e még (korábban package-private, most publikus a HUD és ActionController miatt).}
    \classitem{-recentEvents}{A legutóbbi UI eseményeket tároló adatstruktúra (maximum 12 esemény, FIFO).}
\end{class-template-attribute}
\begin{class-template-method}
    \classitem{+addListener(l : GameStateListener) : void}{Feliratkoztat egy új nézetet az állapotváltozásokra.}
    \classitem{+removeListener(l : GameStateListener) : void}{Eltávolít egy feliratkozott nézetet a listából.}
    \classitem{+fireStateChanged() : void}{Hibatűrő módon (kivételkezeléssel) push-értesítést küld minden feliratkozott listenernek a változásokról.}
    \classitem{+setSelectedName(name : String) : void}{Beállítja a kiválasztott entitást, és automatikusan meghívja a fireStateChanged() metódust.}
    \classitem{+getAllUtak() : Collection}{Publikus getter, amely a GUI számára olvashatóvá teszi az utakat.}
    \classitem{+getAllEntities() : Collection}{Publikus getter az összes entitás lekérdezéséhez.}
    \classitem{+getDepotNode() : String}{Visszaadja a telephely csomópontjának azonosítóját.}
    \classitem{+shortestPath(from : String, to : String) : List}{BFS alapú útvonalkereső, visszaadja a legrövidebb utat két csomópont között.}
    \classitem{+nextStepToward(from : String, to : String) : String}{A BFS algoritmus alapján meghatározza a következő lépéshez szükséges csomópontot a cél felé.}
    \classitem{+roadBetween(nodeA : String, nodeB : String) : Ut}{Visszaadja a megadott két csomópontot összekötő utat.}
    \classitem{+switchActivePlayer() : void}{Átvált a következő játékosra a Pass'N'Play logikának megfelelően.}
    \classitem{+resetActivePlayer() : void}{Visszaállítja a kezdő játékost aktívnak.}
    \classitem{+canControl(j : Jarmu) : boolean}{Segédmetódus, eldönti, hogy az aktuálisan aktív játékos irányíthatja-e a kapott járművet.}
    \classitem{+getRecentEvents() : Deque}{Publikusan elérhetővé teszi az eseménynapló tartalmát a UI számára.}
    \classitem{+enqueueEvent(msg : String) : void}{Beleteszi az új eseményt a recentEvents tárolóba is (max. 12 elem, FIFO elv alapján).}
\end{class-template-method}

\subsubsection{GameActions}
\begin{class-template-responsibility}
    A játék logikai műveleteit összefogó osztály, amely kiegészült a körök lezárását és az NPC-k automatikus mozgatását végző logikával.
\end{class-template-responsibility}
\begin{class-template-interface}
    -
\end{class-template-interface}
\begin{class-template-baseclass}
    -
\end{class-template-baseclass}
\begin{class-template-attribute}
    -
\end{class-template-attribute}
\begin{class-template-method}
    \classitem{+endRound() : void}{Lezárja a kört: végrehajtja az összes NPC autó npcStep() lépését, minden úton alkalmaz egy egységnyi havazást, növeli a körszámlálót (currentRound), frissíti az időt, kiértékeli a játékvége feltételeket, és végül frissíti a nézeteket a fireStateChanged() hívásával.}
\end{class-template-method}

\subsubsection{Auto}
\begin{class-template-responsibility}
    Az utakon közlekedő NPC autókat reprezentáló osztály, amely megkapta a legrövidebb úton (BFS) történő önálló navigációhoz szükséges célpontokat és döntési logikát.
\end{class-template-responsibility}
\begin{class-template-interface}
    -
\end{class-template-interface}
\begin{class-template-baseclass}
    Jarmu
\end{class-template-baseclass}
\begin{class-template-attribute}
    \classitem{+otthon}{Az NPC autó otthonát jelentő csomópont.}
    \classitem{+munkahely}{Az NPC autó munkahelyét jelentő csomópont.}
    \classitem{-aktualisCel}{Az autó éppen aktuális úticélja (otthon vagy munkahely).}
    \classitem{-utolsoCsomopont}{A legutóbb elhagyott csomópont azonosítója a belső állapottartáshoz.}
\end{class-template-attribute}
\begin{class-template-method}
    \classitem{+setupRoute(otthon : String, munkahely : String, startNode : String) : void}{Inicializálja az autó kezdő paramétereit és útvonalát a létrehozáskor.}
    \classitem{+getAktualisCel() : String}{Visszaadja az autó jelenlegi célállomását.}
    \classitem{+npcStep(state : GameState) : void}{Az NPC autó tesz egy lépést a célja felé a legrövidebb úton. Tartalmazza a hibakezelést: ha a kinézett sáv járhatatlan, az autó a csomóponton várakozik.}
    \classitem{-closerNodeToward(node1 : String, node2 : String, target : String) : String}{Privát segédmetódus, amely egy adott út két végpontja közül kiválasztja a célhoz közelebb esőt.}
\end{class-template-method}

\subsubsection{Busz}
\begin{class-template-responsibility}
    A busz specifikus viselkedését megvalósító osztály. A módosítások a rugalmasabb (RegEx / minta alapú) végállomás-felismerést vezették be a statikus string egyezés helyett.
\end{class-template-responsibility}
\begin{class-template-interface}
    -
\end{class-template-interface}
\begin{class-template-baseclass}
    Jarmu
\end{class-template-baseclass}
\begin{class-template-attribute}
    -
\end{class-template-attribute}
\begin{class-template-method}
    \classitem{-isVegallomasNode(nodeName : String) : boolean}{Statikus segédmetódus, amely ellenőrzi, hogy a megadott csomópont neve „Vegallomas" vagy „Vegallomas_" mintázatú-e.}
    \classitem{+onCelUtElerve(target : String, state : GameState) : void}{Felüldefiniált eseménykezelő. Most már az isVegallomasNode helper alapján bármely „Vegallomas_" csomópontnál sikeresen regisztrálja a busz körét.}
\end{class-template-method}

\subsubsection{Hokotro}
\begin{class-template-responsibility}
    A játékos által irányított hókotrót reprezentáló osztály. A változtatások a fej (Söprő/Jégtörő) közvetlen lekérdezését és beállítását teszik lehetővé a grafikus felület és a kezdeti beállítások számára.
\end{class-template-responsibility}
\begin{class-template-interface}
    -
\end{class-template-interface}
\begin{class-template-baseclass}
    Jarmu
\end{class-template-baseclass}
\begin{class-template-attribute}
    -
\end{class-template-attribute}
\begin{class-template-method}
    \classitem{+getAktivFej() : FejTipus}{Visszaadja a hókotróra jelenleg felszerelt aktív fejet (a ContextPanel és ActionController használja).}
    \classitem{+setAktivFej(fej : FejTipus) : void}{Beállítja a hókotró aktív fejét (a főmenü kezdő beállításához és az új hókotró vásárlásához szükséges).}
\end{class-template-method}

\subsubsection{Jatekos}
\begin{class-template-responsibility}
    A játékost és a hozzá tartozó erőforrásokat (vagyon, raktár) reprezentáló osztály. Új funkciója a raktárkészlet biztonságos expozíciója a nézet réteg felé.
\end{class-template-responsibility}
\begin{class-template-interface}
    -
\end{class-template-interface}
\begin{class-template-baseclass}
    -
\end{class-template-baseclass}
\begin{class-template-attribute}
    -
\end{class-template-attribute}
\begin{class-template-method}
    \classitem{+getFejInventory() : List}{Módosíthatatlan (read-only) nézetet ad vissza a játékos birtokában lévő hókotró fejekről, amelyet a GUI (pl. fejcsere dialógus) olvas.}
\end{class-template-method}

\section{Kapcsolat az alkalmazói rendszerrel}
\clearpage
\diagram{docs/img/uml_sequence/hokotrotakaritasview}{Hókotró takarítás}{16cm}
\diagram{docs/img/uml_sequence/jarmukijelolview}{Jármű kijelölés}{16cm}
\diagram{docs/img/uml_sequence/jarmumozgatview}{Jármű mozgatás}{16cm}
\diagram{docs/img/uml_sequence/jatekinditview}{Játék indítás}{16cm}
\diagram{docs/img/uml_sequence/passNPlayhandoffview}{Hand off}{16cm}
\diagram{docs/img/uml_sequence/ujhokotroviewpng}{Új hókotró}{16cm}
\diagram{docs/img/uml_sequence/ujkorview}{Új kör}{16cm}
